% ============================================================================
% 中文注释：第 3 节 贴膜玻璃：光度、协议、面板（2026-11-03 重写）
% ----------------------------------------------------------------------------
% 这一节的任务：把"贴膜玻璃是独立物理体制"从说法变成数字。
%   3.1 膜的成像效应（P1/P2）：反射只占透射 20.1%；反射 34.5% 在高频；
%       暗处反射扰动比亮处大 64%（23.06 vs 14.07 灰阶）。
%   3.2 采集协议与 GT 的危险（P4）：黑布法改光照；AE 状态必须报。
%   3.3 面板与指标：15 场景、手标 Ω（中位 4.6%，2.9–9.2%）、β/α 两个诊断标量。
% 数据出处（服务器）：
%   probe_tinted_photometry_1103.py（P1/P2 数字）
%   remeasure_alpha_annot_1102.py（∂Ω 梯度）
%   qa_split.csv（Ω 面积分布）
% ⚠️ 注意：Ω 内透射被压暗在 1019 上是 0.987（≈不变），只有 165530 是 0.830，
%      所以文中写"分场景、非系统性"，不要写成普适规律。
% ============================================================================
\section{Tinted Glass: Photometry, Protocol, Panel}
\label{sec:benchmark}

% 中文：3.1 膜做了什么——三个可测的量。
\subsection{What the film does to the two layers}
\label{sec:film}

Consider the sensor measurement through a filmed pane,
\begin{equation}
I(x) \;=\; \tau\,T(x) \;+\; \rho\,R(x),
\label{eq:film}
\end{equation}
where $\tau$ is the (film-attenuated) transmission, $\rho$ the reflection coefficient, $T$
the scene and $R$ the reflected scene. Ordinary SIRR implicitly assumes $\rho R$ is
\emph{bright} and $\tau T$ is \emph{well exposed}. Window film inverts both. We measure
three consequences.

% 中文：P1 暗反射。
% 中文：P1 改为"锐度/色度不对称"——这是与 Real20/SIR² 相比真正有 2 倍差异的量。
\paragraph{P1 --- the two layers are not equally sharp.}
The film is laminated on the scene side and the camera is on the bare side, so the two layers
traverse the film differently: transmitted light passes \emph{through} the film bulk (scattered,
absorbed, dyed), whereas the reflection of the behind-camera environment comes off the bare
\emph{glass surface} and never enters the film at all. The prediction is a layer
asymmetry in sharpness and chroma, and this is what we observe. Measured per pair with
$\mathrm{sharp}(X)=\mathbb{E}|\mathrm{HP}(X)|/\mathbb{E}|X|$ (HP: unit-variance high-pass) and
averaged over each dataset:

\begin{center}
\small
\setlength{\tabcolsep}{4pt}
\begin{tabular}{@{}lcccc@{}}
\toprule
 & Real20 & SIR$^2$-Obj & SIR$^2$-Wild & \textbf{ours (film)} \\
\midrule
$\mathrm{sharp}(T)$ & $0.062$ & $0.054$ & $0.041$ & $\mathbf{0.027}$ \\
$\mathrm{sharp}(R)/\mathrm{sharp}(T)$ & $3.32$ & $2.51$ & $2.58$ & $\mathbf{6.76}$ \\
chroma of $T$ & $0.373$ & $0.322$ & $0.393$ & $\mathbf{0.219}$ \\
\bottomrule
\end{tabular}
\end{center}

The transmitted layer of filmed glass is $2.3\times$ softer than on ordinary glass and its
reflection is $2.0\times$ sharper \emph{relative to it}. Mean brightness, mean reflection
magnitude and the high-frequency \emph{fraction} of the reflection are, by contrast,
indistinguishable from Real20 (all measured identically on the four datasets); the asymmetry is what separates the two
regimes.

\paragraph{P2 --- dark content is hit hardest.}
Sorting the scene pixels of each frame by their ground-truth luminance and averaging the
reflection magnitude $|I-T|$ within each quartile gives

\begin{center}
\small
\setlength{\tabcolsep}{4pt}
\begin{tabular}{@{}lcccc@{}}
\toprule
scene luminance & Q1 (dark) & Q2 & Q3 & Q4 (bright) \\
\midrule
mean $|I-T|$ & $\mathbf{23.1}$ & $25.7$ & $17.7$ & $14.1$ \\
\bottomrule
\end{tabular}
\end{center}

The reflection perturbs the darkest quartile $64\%$ more than the brightest. The film
therefore degrades the two layers \emph{in the same places}: where the transmission is
weakest, the relative contribution of the reflection is largest. This is the opposite of
the usual SIRR setting, where a bright reflected scene sits on top of a well-exposed one.

% 中文：P3 透射层信息量低 + 需要注意的口径。
\paragraph{P3 --- the transmission loses information, unevenly.}
Inside $\Om$ the transmitted level relative to $T$ is $0.987$ on average, i.e.\
\emph{not} systematically darker on our panel; the effect is scene-dependent
(min $0.830$ for the darkest case, max $1.046$). We therefore do not claim a universal
darkening inside $\Om$. What is systematic is the loss of \emph{usable} transmitted
detail: inside $\Om$ our models retain only $0.48$--$0.65$ of the ground-truth structure
(Sec.~\ref{sec:structural}), and the strongly-reflected subset carries so little
information that a plain Gaussian blur of the input outperforms every model we trained
(Sec.~\ref{sec:ceiling}).

% 中文：3.2 采集协议与 GT 危险（P4）。
\subsection{Capture protocol and the ground-truth hazard}
\label{sec:protocol}

All scenes are shot on a phone through an automotive window film, handheld, with exposure
and white balance \emph{locked}. For each scene we record a short clip with reflection and
then cover the glass with dark cloth to obtain $T$. There is \textbf{one $T$ per scene}
and many $I$; inputs are extracted at $30$\,fps and cropped to the glass area, and a pair
is kept only if the phase-correlation shift is $\le1.5$\,px and the per-channel gain
deviation is $\le0.10$.

% 中文：P4：黑布法不只是"去掉反射"，它同时改光照。
\paragraph{P4 --- covering the glass is not a neutral operation.}
Covering the pane removes the reflection \emph{and} changes the light reaching the scene;
on filmed glass, where the transmitted signal is already weak, this perturbation is not
negligible relative to the reflection itself. Two practical consequences follow.
\emph{(a) Exposure must be locked}: with auto-exposure active we measured
$\mathrm{median}(I/T)=0.99$ (the covered frame was \emph{brighter}) and a residual layer
correlated with image brightness at $r=0.32$; after locking, $r=-0.01$. \emph{(b) The
residual photometric error must be re-measured, not assumed}: Sec.~\ref{sec:photometric}
shows it is per-image and cannot be shared across scenes. Photometric alignment and
reflection removal therefore have to be handled jointly, and any SIRR benchmark that
reports only pair counts is under-specified.

% 中文：3.3 面板与指标。
\subsection{Panel and metrics}
\label{sec:panel}

The panel contains $15$ scenes ($1280{\times}832$), all outdoor, all shot through the filmed pane
with the reflected environment visible. For every scene we hand-annotated $\Om$ once (median $4.6\%$
of the image, range $2.9$--$9.2\%$). By construction (Sec.~\ref{sec:protocol}) $\Om$ is the
projection of a \emph{dark} object of the reflected environment: the reflection contributes little
there, so $I$ already equals the transmitted layer and a method should only undo the film. Pixels of
the annotated polygon whose transmission deficit is large ($|I-T|>15$) belong to \emph{strong}
reflection rather than to $\Om$; we keep them as a separate \emph{shade} class instead of folding
them in, and report the annulus around $\Om$ separately in Sec.~\ref{sec:axis1}.

We report PSNR, SSIM, HF-PSNR, gradient cosine and the Laplacian-variance ratio
$\mathrm{LapVar}(\That)/\mathrm{LapVar}(T)$, plus region-restricted PSNR inside and outside
$\Om$, and two diagnostic scalars: the residual-reflection coefficient
$\beta=\langle T-\That,\,I-T\rangle/\lVert I-T\rVert^2$ ($-1$: nothing removed, $0$:
perfect) and the per-image gain
$\alpha=\mathrm{median}_i(\overline{\That_i}/\overline{T_i})$, which exposes models that
merely brighten or darken the frame.
